% =====================================================================
\part{Os resultados, com os números na ponta da língua}
% =====================================================================

A Parte III explicou \emph{como} cada número foi obtido. Esta parte organiza
\emph{quais} são, na ordem em que a narrativa os usa. É a que se relê na
véspera.

\chapter{Quarenta anos em três atos}

\section{Ato I --- a pastagem como herança (1985--2000)}

O primeiro mapa da série já mostra uma fronteira em movimento, e não um ponto de
partida natural. Em 1985, quase um terço de Goiás já é pastagem: herança de uma
ocupação que começa antes do primeiro ano da série e que, por isso, este
trabalho \textbf{registra sem medir}.

E o ato é, em volume, \textbf{o mais destrutivo dos três} --- não o preâmbulo dos
outros dois.

\begin{ancora}[Ato I em números]
Pastagem avança \textbf{3,71 Mha} em quinze anos, a \textbf{0,248 Mha/ano} ---
ritmo que nenhum período posterior alcança. \\
Vegetação natural cede \textbf{4,10 Mha}, a \textbf{0,27 Mha/ano} --- quatro
vezes o ritmo dos dois atos seguintes. \\
Lavoura sai de \textbf{1,17 para 3,00 Mha} ($+1{,}83$); a soja quase sextuplica
(\textbf{0,37 para 2,13 Mha}).
\end{ancora}

\begin{armadilha}[O rótulo do ato engana quem o lê depressa]
``A pastagem como herança'' descreve \textbf{o que domina a paisagem}, e não uma
pausa agrícola, que os dados não mostram. A lavoura não fica parada no Ato~I,
apesar da instabilidade macroeconômica do período. O que se mantém no sudoeste é
a \emph{geografia} da lavoura, e não o seu tamanho: o mapa muda pouco de
endereço enquanto a mancha engrossa.
\end{armadilha}

\section{Ato II --- expansão e intensificação (2001--2019)}

\begin{ancora}[Ato II em números]
A agricultura goiana passa de \textbf{9,3\% para 16,0\%} do território. \\
A pastagem atinge o pico por volta de \textbf{2003}, perto de \textbf{14,8 Mha},
e passa a ceder área. \\
No mapa, o sudoeste (Rio Verde, Jataí, Mineiros) converte-se à lavoura sobre o
pasto.
\end{ancora}

\subsection{O ponto que precisa ser dito com todas as letras}

A leitura natural do rótulo erra o alvo. \textbf{Não é a lavoura que acelera em
2001.}

\begin{ancora}[A comparação que desfaz a leitura fácil]
Lavoura: cresce 2,31 Mha em dezoito anos, a \textbf{0,128 Mha/ano}, contra
\textbf{0,122} no Ato~I --- praticamente o mesmo ritmo absoluto. \\
Soja: avança a \textbf{0,116 Mha/ano}, contra \textbf{0,117} no ato anterior. \\
\textbf{O que inverte o sinal é a pastagem}, que passa de $+0{,}248$ para
$\mathbf{-0{,}076}$ Mha/ano.
\end{ancora}

Daí decorre a leitura correta do ato: o que muda em 2001 não é a velocidade da
expansão agrícola, e sim a \textbf{fonte da terra} que a alimenta. Até então a
lavoura crescia enquanto o pasto \emph{também} crescia, ambos sobre vegetação
natural; a partir dali a lavoura passa a crescer \emph{sobre o pasto}.

\begin{nabanca}
``A quebra que o método primário detecta em 2001 é, na origem, uma quebra da
série de \emph{pastagem}, e é por isso que ela aparece num teste aplicado
conjuntamente às três séries, e não numa delas isolada. É o mesmo mecanismo que
a Perna 2 mede como substituição local e que a Perna 4 reencontra, vinte anos
depois, na freada do Sul.''
\end{nabanca}

\subsection{O contexto, que entra como cenário e não como achado}

Convém ser explícito sobre isto, porque é um convite a excesso de afirmação. A
entrada da China na OMC no fim de 2001, a sistematização do crédito rural a
partir de 2002--2003 e o super-ciclo de preços de \emph{commodities} entram
\textbf{para situar o leitor no período}, e a eles não se pendura conclusão
alguma. Dos três, apenas o super-ciclo chegou a entrar em teste --- no desenho
de diferenças em diferenças que foi rebaixado por não existir grupo não tratado
--- e nenhum dos três é sustentado por fonte conferida.

\section{Ato III --- conversão acelerada sob rótulo ambíguo (2020--2024)}

Este é o ato em que o satélite engana, e ele já foi destrinchado no
Capítulo~\ref{cap:markov} e antes. Os números da contradição, para tê-los
juntos:

\begin{ancora}[A contradição e a sua resolução]
Pela classe ``Agricultura'': a conversão de pasto em lavoura no Sul goiano
\textbf{cai 88\%}. \\
Pela soja que o IBGE levanta em campo: a área \textbf{cresce 38\%} no estado na
mesma janela. \\
Pela união de ``Agricultura'' com ``Mosaico'', imune à reetiquetagem por
construção: a taxa no Sul \textbf{acelera 51\%}. \\[4pt]
\textbf{O que parou foi a classe, e não o território.}
\end{ancora}

\begin{ancora}[O resto do ato]
O recuo da pastagem quase quadruplica de velocidade: de \textbf{0,07 para 0,27
Mha/ano}. \\
A vegetação natural \emph{não} muda de ritmo: perde cerca de \textbf{0,06
Mha/ano} tanto no Ato~II quanto no Ato~III. O que muda é que ela passa a ceder
\textbf{ao norte}. \\
Goiás chega a 2024 com \textbf{34,9\%} do território em vegetação natural, sem
sinal de estabilização até aqui.
\end{ancora}

\begin{nabanca}
``O que mudou no Ato~III não foi a velocidade da perda de vegetação, e sim o seu
endereço.''
\end{nabanca}

\section{O balanço de quarenta anos}

\begin{ancora}[A tabela que precisa estar de cor]
Vegetação natural: \textbf{17,65 $\to$ 11,88 Mha} ($-5{,}8$), de \textbf{51,9\%
a 34,9\%} do território. \\
Agricultura: \textbf{1,17 $\to$ 5,58 Mha} ($\times 4{,}8$). A soja cresce
$\times 12$ em área (0,37 $\to$ 4,50 Mha) e $\times 13$ em produção
(1,3 $\to$ 17,0 Mt). \\
Pastagem: \textbf{10,98 $\to$ 11,99 Mha} ($+1{,}0$) --- saldo enganoso, porque a
trajetória é em U invertido, com pico próximo de 14,8 Mha em 2003. \\
Lotação bovina: \textbf{1,45 $\to$ 1,94 cab/ha}. O rebanho cresce \textbf{46\%}
(15,9 $\to$ 23,2 milhões de cabeças) e a área de pasto apenas \textbf{9\%}: a
pecuária se intensifica. \\
Mosaico de Usos: \textbf{3,63 $\to$ 3,59 Mha} ($\approx 0$) --- saldo nulo,
caminho não: cai de 10,7\% do estado a 6,1\% em 2019 e volta a 10,5\% em 2024.
\end{ancora}

\begin{armadilha}[Por que os fluxos não fecham por subtração com a tabela]
Se a banca tentar conciliar os 4,65 Mha de entrada em agricultura com o
acréscimo de 4,41 Mha da tabela, a resposta tem \emph{três} partes: as duas
convenções de classe (campo alagado, 0,70 Mha; silvicultura, 0,16 Mha) e o fato
de que o cruzamento entre as duas pontas conta o estado inicial e o final do
pixel, e não cada passagem intermediária.
\end{armadilha}

\chapter{Frente 1 --- o padrão}

\section{O que ficou estabelecido}

A conversão e o rebanho deslocam-se ao norte, e o gradiente latitudinal
persiste. É a única das quatro frentes classificada como \textbf{estabelecida}.

Os números do deslocamento estão no Capítulo~\ref{cap:centro}. O que falta aqui
é o \emph{ritmo}, o \emph{vão} e as verificações que não são de método.

\section{O ritmo, que diz mais que o total}

\begin{ancora}[Quilômetros por ano, por ato]
Pastagem: $+2{,}5$, $+1{,}4$, $\mathbf{+2{,}8}$ \\
Rebanho: $+1{,}8$, $+1{,}4$, $+2{,}0$ \\
Agricultura: $+2{,}0$, $+1{,}9$, $+0{,}1$
\end{ancora}

Pastagem e rebanho desaceleram no Ato~II e \textbf{voltam a acelerar} no
Ato~III, quando a pastagem marcha mais rápido do que em qualquer momento da
série.

\begin{quote}\itshape\color{cortinta}
O achado desta frente não é que a marcha existiu e terminou, e sim que ela chega
ao fim da série no seu ritmo mais alto.
\end{quote}

A queda da agricultura para $+0{,}1$ no Ato~III é o único dos nove números que
não é fenômeno de território: é a linha achatando pela mudança de rótulo.

\section{A componente leste, que a palavra ``norte'' não carrega}

Reportar só a componente norte seria escolher a projeção que convém. As duas são
calculadas e as duas recebem o mesmo tratamento de incerteza.

\begin{ancora}[Os azimutes]
Pastagem: $+19$ km a leste contra $+78$ ao norte $\Rightarrow$ azimute
\textbf{14°}. \\
Rebanho: $+23$ contra $+67$ $\Rightarrow$ azimute \textbf{19°}. Quase
meridionais. \\
Agricultura: $+49{,}5$ km a leste, IC $[+22{,}6; +77{,}1]$, contra $+65{,}2$ ao
norte $\Rightarrow$ resultante de \textbf{82 km} em azimute de \textbf{37°} ---
uma marcha ao \emph{nor-nordeste}.
\end{ancora}

\begin{naodiz}
A componente leste é \textbf{menos robusta} que a norte --- nenhuma delas
atravessa a varredura de blocos: a da pastagem inclui zero já em blocos de três
AMCs, a da agricultura em oito, a do rebanho em catorze. Por isso a marcha continua
sendo descrita pelo eixo meridional. O que ela \emph{sugere}, e este trabalho
não testa, é a aproximação da lavoura goiana ao eixo do Leste Goiano e do Vão do
Paranã, que é por onde Goiás faz divisa com o Cerrado baiano.
\end{naodiz}

\section{O vão que nunca se fecha}

\begin{ancora}[Os ``dois Goiáses'' em métrica]
Em todos os quarenta anos, a lavoura permanece entre \textbf{123 e 135 km} ao
sul do pasto e do rebanho. \\
O vão nunca se fecha e nunca ultrapassa essa faixa, mas não é constante:
estreita-se até 123 km em 2018--2019 e volta a 135 km em 2024. \\[4pt]
Em 2024, medido em três réguas: \textbf{135 km} pela classe do satélite,
\textbf{111 km} pela união, \textbf{122 km} pela soja que o IBGE mede em campo.
Grande e persistente em qualquer uma das três.
\end{ancora}

\begin{armadilha}
Não leia a reabertura recente (123 $\to$ 135 km) como afastamento real: ela é
medida com o centroide agrícola já contaminado pela mudança de rótulo.
\end{armadilha}

\section{A abertura por formação, que corrige a leitura do agregado}

Autocorreção que vale apresentar, porque ela mostra que o agregado escondia
comportamentos distintos.

\begin{ancora}[A vegetação natural aberta]
Formação florestal: $+8{,}7$ km, IC $[+2{,}5; +15{,}1]$ --- \textbf{andou}. \\
Campo nativo: direção robusta, magnitude imprecisa (intervalo largo) ---
\textbf{andou}. \\
Formação savânica, a mais extensa das três e por isso a que domina o agregado:
intervalo \textbf{inclui zero}.
\end{ancora}

\begin{nabanca}
``O que parecia imobilidade da vegetação é a soma de uma classe parada com duas
que andaram. O número médio do agregado escondia comportamentos distintos por
dentro.''
\end{nabanca}

\section{Os dois controles negativos}

São o que mostra que não se trata de deslocamento genérico --- de ``o mapa
inteiro subindo''.

\begin{ancora}[Os controles]
\textbf{Área urbana:} anda \textbf{8 km para o sul}, IC $[-19; -1{,}5]$ --- no
sentido \emph{oposto} ao da fronteira. \\
\textbf{Leite}, pecuária da bacia consolidada do Sul: sobe \textbf{$+30$ km},
menos da metade do que sobe o rebanho total ($+67$). E o vão entre os dois
centros \textbf{mais que dobra}: de 30 km em 1985 para 68 km em 2024.
\end{ancora}

\begin{quote}\itshape\color{cortinta}
Quem marcha é a pecuária de corte na fronteira, e não o mapa inteiro subindo.
\end{quote}

\section{A camada do fogo, e o cuidado que ela exige}

\begin{ancora}[O fogo como vanguarda geográfica]
O centro de massa do fogo em vegetação natural fica ao norte do centro da
conversão em \textbf{todos os 39 anos} em que as duas séries coexistem
(1985--2023), com dianteira média de \textbf{73 km}, \emph{sem nunca trocar de
lado}.
\end{ancora}

\begin{naodiz}[E isso é tudo o que ele é]
A leitura de que o fogo abriria caminho foi \textbf{testada e não se sustenta no
tempo}: a liderança temporal desaparece sob transformação logarítmica, inverte
quando se removem os anos de seca de 1985 e 2010, e o teste de precedência
agregado é nulo. \\[4pt]
\textbf{O fogo marca onde a fronteira está, não quando ela vai andar.}
\end{naodiz}

\begin{armadilha}[Detalhe de contagem]
São 39 anos e não 40 porque a série de conversão termina um ano antes: cada
transição é rotulada pelo ano de \emph{origem} do par.
\end{armadilha}

\chapter{Frente 2 --- o mecanismo}

\section{As duas lógicas}

No Sul, a transição dominante é de pastagem para lavoura: intensificação, terra
que troca de função. No Norte, é de vegetação natural para pastagem: fronteira,
terra nova sendo aberta.

A medida que sustenta a divisão é a que \textbf{não depende da classe ambígua},
já que nem sua origem nem seu destino passam por ela:

\begin{ancora}[A assimetria entre as pontas do estado]
A conversão de vegetação natural em pastagem cai \textbf{49\% no Sul Goiano}
entre os Atos~II e III e apenas \textbf{13\% no Norte} --- na malha das cinco
mesorregiões. Quase quatro para um.
\end{ancora}

\section{O que a idade da pastagem estabeleceu, e o que ela não estabeleceu}

Os números da mistura estão no Capítulo~\ref{cap:mistura}. O saldo desta
subfrente é uma \textbf{tripla negativa} que vale saber enunciar sem soar
derrotado.

\vspace{4pt}
\textbf{Não é a região.} A mesorregião explica \textbf{0,5\%} da variação da
idade.

\textbf{Não é o sistema produtivo.} Para o plantio direto, a correlação simples
confirma a previsão, mas o indicador também se concentra no Sul; segurando o
gradiente espacial, a associação encolhe cerca de um terço, fica na fronteira da
significância sem cruzá-la e não sobrevive à correção para testes múltiplos. O
veredito é \textbf{não estabelecido}, e a distinção importa, porque um nulo
limpo encerraria a pergunta, ao passo que um indeterminado não encerra.

\textbf{E o crédito aponta na direção contrária à esperada.} A associação é a
única que resiste ao controle espacial, à correção para testes múltiplos e à
troca de régua, e na régua imune ela \emph{se fortalece}, de $+0{,}22$ para
$+0{,}30$. Mas o sinal diz que \emph{onde o crédito cresceu mais, o pasto
convertido é mais velho}, o que é compatível com reserva sendo ativada por
dinheiro novo, e não com crédito financiando rotação.

\begin{naodiz}[As duas ressalvas do resultado do crédito]
Ele só existe quando os anos recentes entram: no recorte até 2019 é
praticamente nulo. E o mecanismo não foi investigado.
\end{naodiz}

\begin{nabanca}[O achado negativo mais útil do trabalho]
``A escolha entre girar o capim em três anos ou deixá-lo trinta acontece abaixo
da escala que estes dados alcançam. Ela não se lê na região, não se lê no
sistema produtivo do município e mal se lê no crédito que chegou até ele. A
consequência é que a próxima pergunta desta linha exige outra \emph{unidade de
observação} --- dado por imóvel rural --- e não mais tratamento estatístico do
mesmo agregado.''
\end{nabanca}

\begin{armadilha}[Sobre o plantio direto]
O Censo Agropecuário não tem variável de integração lavoura-pecuária. O
indicador mais próximo disponível é o plantio direto, que é \textbf{prática de
conservação de solo, e não de integração com pecuária} --- ainda que ambas
costumem andar juntas na lavoura tecnificada do Cerrado. Diga isso antes de a
banca dizer.
\end{armadilha}

\chapter{Frente 3 --- o empurrão e o motor comum}

\section{Por que a hipótese do empurrão era irresistível}

As duas frentes anteriores entregam dois fatos que, postos lado a lado, quase
escrevem sozinhos uma conclusão: a fronteira inteira subiu ao norte, e subiu
fazendo coisas opostas em cada metade do estado. Somados, sugerem que o pasto
expulso do Sul foi reaparecer no Norte, sobre o Cerrado.

\begin{quote}\itshape\color{cortinta}
É a hipótese que mais favorecia a tese inicial deste trabalho. Ela tem, porém,
um problema de lógica antes de ter um problema de dados: uma coincidência no
espaço não é um mecanismo, e dois processos independentes movidos pela mesma
força externa desenham o mesmo mapa que um empurrão desenharia.
\end{quote}

Separar as duas leituras exige procurar uma assinatura que \emph{só} o empurrão
produz. Foram duas, e ambas específicas o bastante para poder falhar --- o que
as torna teste, e não ilustração.

\section{Os dois resultados, em números}

\begin{ancora}[No espaço: doze especificações]
$3$ réguas de lavoura $\times$ $2$ janelas $\times$ $2$ desfechos $=$
\textbf{12}. \\
\textbf{Com oito vizinhos, as doze estimativas são negativas.} A hipótese
exigia positivas. Com quatro e com doze vizinhos são 10 e 8 de 12, e nenhuma
positiva é significativa. \\
Onze não são significativas; a décima segunda tem $p = 0{,}02$ e é negativa ---
e está na régua contaminada pelo Mosaico.
\end{ancora}

\begin{ancora}[No tempo: vinte e quatro combinações]
$p$ mínimo $= \mathbf{0{,}078}$; nenhuma combinação é significativa. \\
Poder: \textbf{93\%} contra um empurrão forte, \textbf{48\%} contra um moderado.
\end{ancora}

Somadas, são \textbf{36 recortes}, e é esse o número que a Discussão usa.

\section{O que ocupa o lugar do empurrão}

A \textbf{substituição local} é o maior efeito de toda a frente, com coeficiente
de $-0{,}52$ a $-1{,}14$ nas células medidas em área de satélite. Os detalhes de
leitura estão no Capítulo~\ref{cap:espacial}.

\begin{nabanca}[O posicionamento na literatura, que vale ouro aqui]
``É o desfecho observado da disputa por terra que a literatura de intensificação
projeta em modelo: no Sul goiano ela aparece resolvida, e resolvida a favor da
lavoura. A literatura de intensificação e a de deslocamento costumam ser
mobilizadas em conjunto, como se o mesmo dado sustentasse as duas; aqui elas se
separam empiricamente. O mecanismo local aparece em quase toda especificação, ao
passo que o mecanismo à distância não aparece em nenhuma.''
\end{nabanca}

\section{O motor comum, em uma frase e três ressalvas}

O encadeamento suposto é direto e vale dizê-lo assim: Goiás produz soja e carne
para exportação, ambas cotadas em dólar; quando o real se desvaloriza, a mesma
tonelada exportada rende mais reais ao produtor sem que o preço internacional
mude; a rentabilidade de plantar e de criar sobe, e a procura por terra cresce.

\begin{quote}\itshape\color{cortinta}
O câmbio não cria fazenda; torna compensador avançar.
\end{quote}

Os números e as ressalvas estão no Capítulo~\ref{cap:shiftshare}. O que importa
memorizar é o veredito: \textbf{corroborante, não estabelecido}, e a razão de
ele não decorrer do nulo anterior.

\begin{armadilha}[Uma armadilha lógica que a banca pode armar]
``Você descartou o empurrão, logo o motor comum está estabelecido.'' Não. A
hipótese do \emph{drive} comum \textbf{não decorre} de o empurrão não ter
aparecido, porque descartar uma explicação não instala a seguinte. Ela foi
testada por conta própria, e o resultado dela é o que é.
\end{armadilha}

\section{A precedência de infraestrutura e crédito}

Verificação independente, e de raciocínio limpo: quem puxa costuma estar à
frente. Se silos, bancos ou lavouras puxassem a marcha, seus centros de gravidade
estariam \emph{ao norte} do pasto e do rebanho; se apenas consolidassem o que já
foi aberto, estariam ao sul.

\begin{ancora}[Nenhuma camada aparece à frente]
Crédito rural: cerca de \textbf{70 km ao sul} da pastagem, e o vão não se
fecha, pois era 77 km em 2013 e segue em 69 km em 2024. \\
Capacidade de armazenagem: cerca de \textbf{150 km ao sul} do pasto,
praticamente sobre o núcleo de lavoura do Sudoeste, a \textbf{16 km} dele. \\
Cadeia exportadora: co-move sem liderar.
\end{ancora}

\begin{naodiz}
O resultado é \textbf{descritivo}: diz onde cada camada está, e não que ela seja
incapaz de puxar. É posição, não efeito. O que ele faz é contrariar a suposição,
comum no debate, de que restringir infraestrutura ou crédito seja o ponto de
aplicação natural para deter uma fronteira que já se moveu.
\end{naodiz}

\chapter{Frente 4 --- o teto de oferta}

\section{As duas precisões que abrem a frente}

\textbf{Primeira: não foi a lavoura que freou.} Essa leitura seria o artefato de
rótulo já desarmado. O que freou é a abertura de vegetação natural, medida pela
conversão de vegetação em pastagem --- transição cuja origem e cujo destino
ficam fora da classe ambígua.

\textbf{Segunda: não foi o estado que desacelerou, e sim o Sul.}

\begin{ancora}[O agregado não se moveu; o endereço, sim]
Fluxo bruto de conversão de vegetação natural em Goiás: \textbf{0,071 Mha/ano}
no Ato~II para \textbf{0,072} no Ato~III --- praticamente o mesmo número. \\
O que mudou foi o endereço: o fluxo bruto \textbf{caiu 37\% na região Sul} e
subiu no Centro e no Norte (malha de três regiões).
\end{ancora}

\begin{quote}\itshape\color{cortinta}
A marcha não perdeu velocidade; mudou de lugar. A pergunta desta frente não é
por que a fronteira parou --- ela não parou ---, e sim o que a expulsou do Sul.
\end{quote}

\begin{armadilha}[Dois números de queda que não são comparáveis]
Os \textbf{49\%} do Sul Goiano e os \textbf{37\%} da região Sul são medidas
diferentes em malhas diferentes: cinco mesorregiões contra três regiões. Não os
apresente como o mesmo número arredondado de dois jeitos.
\end{armadilha}

\section{A demanda estava no pico}

\begin{ancora}[Os três sinais de demanda, médias dos dois períodos]
Índice de câmbio real efetivo: \textbf{134,5 $\to$ 169,0} \\
Índice de preço recebido pela soja: \textbf{104,4 $\to$ 186,4} \\
Crédito rural de Goiás: \textbf{14,3 $\to$ 24,1} bilhões de reais de dez./2024 \\
Soja efetivamente plantada: \textbf{$+38\%$} na mesma janela
\end{ancora}

\begin{armadilha}[A média do crédito esconde o formato da série]
Ele dispara em 2021--2022 (\textbf{29,6} e \textbf{34,2} bilhões) e desaba em
2023--2024, terminando em \textbf{15,8} --- abaixo da média de 2013--2019.
Câmbio e preço permanecem altos até o fim da janela; o crédito se inverte no
último biênio. São eles e a área de soja plantada, \textbf{e não a média do
crédito}, que sustentam a leitura de demanda alta no Ato~III. Diga isso antes de
a banca encontrar.
\end{armadilha}

\section{O teste que fecha a questão não é o do estado, e sim o do próprio Sul}

Se o Sul tivesse perdido apetite por terra, teria parado de converter
\emph{qualquer} terra.

\begin{ancora}[O Sul não parou de querer terra]
No Sul do Ato~III, a taxa anual de conversão de pastagem em lavoura ou uso misto
\textbf{sobe 51\%} em relação ao Ato~II. \\
O ritmo de expansão da soja plantada --- os hectares que o IBGE registra como
acréscimo a cada ano --- \textbf{sobe 244\%}. \\
Ao mesmo tempo, a abertura de Cerrado novo \textbf{cai pela metade}.
\end{ancora}

\begin{armadilha}
Os 38\% e os 244\% \textbf{não são a mesma medida}: 38\% é o crescimento do
estoque plantado no \emph{estado} entre as pontas da janela; 244\% é a
aceleração do \emph{acréscimo anual} no \emph{Sul} entre um ato e outro.
\end{armadilha}

\begin{quote}\itshape\color{cortinta}
O Sul não parou de querer terra. Parou de tirá-la do Cerrado e passou a tirá-la
de pasto já aberto.
\end{quote}

\section{O estoque ainda exposto}

\begin{ancora}[O número que a justificativa do trabalho usa]
Restam \textbf{6,56 Mha} de vegetação natural ainda exposta à conversão em
Goiás --- cerca de \textbf{60\%} do que havia em 1985 --- e o \textbf{Norte
concentra 44\%} deles. \\
Em 1985, o mesmo estoque era de \textbf{10,67 Mha}; foram suprimidos
\textbf{4,11 Mha}.
\end{ancora}

\subsection{Por que ``convertível'' é savana e campo, e não toda a vegetação}

A definição não é arbitrária: decorre do que o próprio esgotamento mostra.

\begin{ancora}[Perdas por formação, 1985--2024]
Formação savânica: $-38{,}2\%$ \\
Campo nativo: $-42{,}5\%$ \\
Formação florestal: $-22{,}8\%$ \\[4pt]
Savana e campo são as classes que a fronteira \emph{de fato} consome.
\end{ancora}

\begin{naodiz}[O vocabulário é obrigatório aqui]
O termo ``convertível'' mede \textbf{exposição, e não disponibilidade}. É o
Cerrado que ainda \emph{pode ser perdido}, e não uma reserva de terra que este
trabalho trate como aproveitável. Tratá-lo como reserva inverteria o sentido do
resultado. \\[4pt]
E é um \textbf{teto}: sem cadastro ambiental integrado pixel a pixel, o estoque
inclui Reserva Legal e APP que na prática não podem ser convertidas. A terra
realmente exposta é \emph{menor} que 6,56 Mha, nunca maior. Um cenário
descontando 20\% de Reserva Legal foi executado e não altera o ordenamento entre
regiões.
\end{naodiz}

\section{Menos terra, ou terra pior?}

A alternativa mais séria à leitura de esgotamento, e o trabalho a testa em dois
canais que dão respostas diferentes.

\textbf{O mecanismo existe:} a propensão a converter é maior onde a aptidão é
maior ($\beta = +0{,}004$ no Sul; $p < 0{,}01$), nas três regiões e em todos os
cortes de estoque testados. A pergunta é se ele foi longe o bastante para mudar
a natureza do que sobrou.

\begin{ancora}[Entre unidades: não foi]
Aptidão média do estoque ainda exposto, ponderada pelo próprio estoque, no Sul:
\textbf{4,60} em 1985 e \textbf{4,61} em 2024. \\
A fronteira \emph{não} recuou para municípios ruins. \\[4pt]
Curiosamente, quem exibe o efeito é o \textbf{Norte}, cujo estoque perde
\textbf{0,12} ponto de aptidão no período --- lá, e não no Sul, a conversão está
de fato consumindo primeiro as unidades melhores, que é o comportamento de uma
fronteira ativa descendo o gradiente.
\end{ancora}

\begin{ancora}[Dentro da unidade: o quadro é outro]
Fração do remanescente em formação florestal (mata de galeria e cerradão): \\
Sul: \textbf{52,2\% $\to$ 59,9\%} \\
Norte: \textbf{29,6\% $\to$ 34,3\%} \\
Estado: \textbf{37,0\% $\to$ 42,5\%}
\end{ancora}

\begin{naodiz}[A leitura exata, que é mista e não deve ser arredondada]
A alta é do estado inteiro, e o excesso do Sul sobre a média estadual não passa
de dois pontos. A \emph{tendência} \textbf{não separa} seleção de uma oscilação
do classificador na borda entre savana e formação florestal --- resíduo que este
trabalho não mediu. \\[4pt]
O que distingue o Sul é o \textbf{nível}, e o nível não depende da tendência: em
2024, \textbf{três quintos} de tudo o que o Sul ainda tem de vegetação natural
são a fisionomia que a fronteira não consome, contra um terço no Norte. E o Sul
já entra na série com 52\%, de modo que a diferença é de \textbf{fisiografia
antes de ser de história}. \\[4pt]
Portanto: a alternativa de qualidade está \textbf{eliminada} no canal
\emph{entre} unidades e \textbf{viva} no canal \emph{dentro} delas --- viva como
\emph{estado}, e não necessariamente como processo.
\end{naodiz}

\begin{nabanca}
``Ela não substitui a leitura de esgotamento, porque o Sul de fato tem menos.
Mostra que `acabou a terra convertível' e `o que restou é mata de galeria em
relevo quebrado' descrevem, no Sul, largamente o mesmo estado de coisas. E é
exatamente por isso que a minha definição de estoque convertível exclui a
floresta.''
\end{nabanca}

\section{A proteção integral não estava no caminho}

Uma fronteira pode parar por duas razões muito diferentes --- a terra acabou, ou
a lei a barrou --- e as duas produzem o mesmo gráfico enquanto pedem políticas
opostas. Para que o teto fosse institucional, a proteção teria de dar conta de
duas coisas: o \emph{momento} em que o Sul freou e o \emph{lugar} para onde a
fronteira foi. Não dá conta de nenhuma.

\begin{ancora}[No tempo: ela não se moveu]
\textbf{76\%} da proteção integral que Goiás tem hoje já existia em 1985; \\
\textbf{89\%} já existia em 2000. \\
Em quarenta anos ela cresceu \textbf{0,12 Mha}, contra \textbf{4,11 Mha}
suprimidos do mesmo estoque convertível. \\
Uma variável praticamente constante não data uma inflexão.
\end{ancora}

\begin{ancora}[No espaço: ela não está à frente]
Dos 6,56 Mha ainda expostos, \textbf{6,35 estão fora} de Unidade de Conservação
de Proteção Integral: \textbf{97\%} pela malha de unidades de conservação e
\textbf{94,3\%} quando a mesma conta é refeita pixel a pixel, que é a régua mais
exigente. \\
A Proteção Integral cobre \textbf{menos de 3\%} em qualquer região.
\end{ancora}

\begin{naodiz}[E aqui o vocabulário precisa ser cirúrgico]
``Desprotegido'', no sentido produzido por esse recorte, é \textbf{estreito}:
significa fora do grupo cujo regime veda a conversão de modo geral. \\[4pt]
\textbf{Não} significa livre para desmatar. \textbf{Não} significa que nada no
grupo de Uso Sustentável restrinja conversão --- há categorias nele que
restringem, e o trabalho não as examinou unidade a unidade. \\[4pt]
Significa apenas que \emph{a única barreira de que este trabalho tem medida
espacial completa não está no caminho da fronteira}. Reserva Legal e APP seguem
valendo dentro de cada propriedade e protegem parte desse mesmo estoque, e é
por isso que os 6,56 Mha são um teto.
\end{naodiz}

\begin{armadilha}[O limite que fica em aberto, e que é honesto declarar]
Reserva Legal e APP incidem dentro de cada propriedade, ganharam fiscalização
efetiva com o Cadastro Ambiental Rural a partir de 2012 --- \emph{na janela de
interesse} --- e tendem a morder mais forte onde o remanescente já é escasso. No
Sul esgotado, ``acabou a terra convertível'' e ``o que restou é legalmente
inconversível'' produzem a mesma série observada. \\[4pt]
O cenário de $-20\%$ compara regiões entre si e \textbf{não separa} essas duas
leituras dentro do Sul. Separá-las exigiria o cadastro ambiental integrado pixel
a pixel, que este trabalho não teve, e é a primeira das arestas do
cronograma.
\end{armadilha}

\chapter{As duas contas que fecham o trabalho}

\section{Carbono}

A aritmética e as decisões D30 e D31 estão no Capítulo~\ref{cap:carbono}. Aqui,
os números da narrativa.

\begin{ancora}[Quem paga]
Formação savânica: \textbf{573 Mt} de estoque removido. \\
Formação florestal: \textbf{340 Mt}. \\[4pt]
\textbf{Não é um resultado sobre densidade, e sim sobre área:} a savânica perdeu
2,6 vezes mais superfície, e isso mais que compensa o fato de valer menos por
hectare. O desconto do uso entrante não altera o ordenamento: a razão passa de
1,69 para 1,58.
\end{ancora}

\begin{ancora}[Mas o custo não é proporcional à área]
A formação florestal responde por \textbf{um quarto} da área perdida e por
\textbf{mais de um terço} do carbono. \\
Um hectare dela vale \textbf{1,6} hectare de savânica e \textbf{2,6} hectares de
formação campestre. \\
E o desconto do uso entrante \emph{reforça} essa afirmação: subtrair um mesmo
valor de dois estoques desiguais aumenta a razão entre eles, e o hectare de
floresta passa a valer \textbf{1,7} de savânica.
\end{ancora}

\begin{ancora}[Quando se pagou]
Três quartos (\textbf{722 Mt}) saíram no \textbf{Ato~I}, a \textbf{48,1 Mt/ano},
com o Sul à frente no período (275 Mt, ou 38\% do ato, a 18,3 Mt/ano). \\
Depois de 2001 o ritmo cai cerca de \textbf{cinco vezes}, para \textbf{9,7
Mt/ano} no Ato~II. \\
E a geografia se inverte: no Ato~III o Sul é o que \emph{menos} remove (1,9
Mt/ano) e o Centro o que mais (5,3). \\
O centroide da perda marcha \textbf{91 km ao norte}, atrás da fronteira. \\[4pt]
No acumulado dos quarenta anos, as três regiões respondem por parcelas próximas
(Sul 34\%, Norte 34\%, Centro 32\%): \textbf{o que distingue o Sul é ter
removido cedo}, e não o total acumulado.
\end{ancora}

\begin{ancora}[A composição, que fecha a conta]
No Ato~I, o estoque removido se dividia quase meio a meio entre as duas
formações lenhosas: \textbf{45\%} de mata de galeria e cerradão contra
\textbf{49\%} de savânica. \\
No Ato~III, \textbf{93\%} vem da savânica e apenas \textbf{1\%} de floresta.
\end{ancora}

\begin{nabanca}
``A fronteira deixou de consumir a formação mais densa em carbono e passou a
consumir quase só a mais extensa. Isso baixa o custo climático de cada hectare
convertido ao mesmo tempo que o desloca para a classe de que o Cerrado
remanescente é feito.''
\end{nabanca}

\begin{armadilha}[Duas armadilhas de contagem]
\textbf{Primeira:} os três atos somam \textbf{944 Mt} e não os 973 do período
inteiro, porque cada ato é medido entre o seu primeiro e o seu último ano e as
duas viradas não pertencem a nenhum deles. A diferença é de 3\%. \\[4pt]
\textbf{Segunda:} o ritmo não segue caindo depois do Ato~II. Somadas as três
regiões, o Ato~III remove \textbf{11,8 Mt/ano}, \emph{acima} do Ato~II --- e
essa comparação depende da convenção de denominador (intervalos ou anos), por
isso não é reportada como percentual. \\[4pt]
E nos parágrafos por ato a palavra exata é \textbf{``removido''}, e não
``emitido''.
\end{armadilha}

\section{Desenvolvimento municipal}

A conta mais delicada do trabalho, e a que exige a maior contenção.

\begin{ancora}[Os números]
Entre 2013 e 2021, a fronteira norte quase dobrou a área cultivada:
\textbf{$+93\%$}, contra \textbf{$+14\%$} no Sul. \\
O IFDM subiu em toda parte no período (em Goiás, de \textbf{0,49 a 0,63}). \\
Mas subiu no Norte na \emph{mesma medida} que no Sul, e o vão permanece intacto:
\textbf{$-0{,}083$} em 2023, IC 95\% $[-0{,}108; -0{,}058]$, contra
\textbf{$-0{,}092$} em 2013. \\[4pt]
Num painel de efeitos fixos dentro do município, \textbf{dobrar a área vale
0,008 ponto de IFDM}, contra os cerca de \textbf{0,15} que o índice subiu na
década.
\end{ancora}

\begin{naodiz}[As quatro ressalvas, todas obrigatórias]
A leitura é \textbf{associativa}, com o gradiente espacial controlado, e não
causal. \\
O IFDM é \textbf{índice composto e municipal}, e não capta distribuição interna. \\
A janela é a da \emph{série revista} da FIRJAN, que começa em 2013 e \textbf{não
é emendável} com a anterior --- onze anos é janela curta para um índice de
desenvolvimento. \\
Estabelecer o contrário exigiria outro desenho.
\end{naodiz}

\begin{nabanca}
``Com essas ressalvas, o que se pode dizer é que a expansão da área cultivada,
no período e na escala observados, não aparece associada a convergência de
desenvolvimento. Isso recomenda cautela com o argumento de que abrir área é, por
si, política de desenvolvimento regional. Não é uma afirmação sobre o que a
expansão \emph{causa}.''
\end{nabanca}
